\paragraph{Result R63: all holonomies reconstruct the original source.}
The complete proof (H1)--(H66), including (H49a)--(H49g), and the
typed continuation (HDI1)--(HDI24) retain the original half-density
source $\Psi_N=\mathcal U_k\mathcal V|_{\mathcal P_N}$, the full
relative Hilbert fibre, and $g=2\xi$. For $L>0$ the exact map and
inverse are
\[
\begin{gathered}
 \mathcal Z_{L,\theta}\psi(r)=
       \sum_{a\in\mathbb Z}e^{ia\theta}\psi(r+aL),\qquad
 \psi(r+aL)=\int_0^{2\pi}e^{-ia\theta}
       \mathcal Z_{L,\theta}\psi(r)\,\frac{d\theta}{2\pi},\\
 \int_0^{2\pi}M_{N,\theta}\,\frac{d\theta}{2\pi}=M_N,
 \qquad M_{N,\theta}=(\mathcal Z_{L,\theta}\Psi_N)^*
                             (\mathcal Z_{L,\theta}\Psi_N).
\end{gathered}
\]
Fourier series in the original translation index prove the isometry.
The quasi-periodic basis has coefficients
$L^{-1/2}\widehat\psi((2\pi n+\theta)/L)$; the full zero mode
remains. The length-$mL$ source has the exact finite character
decomposition $M_N(mL,0)=m^{-1}\sum_jM_N(L,2\pi j/m)$,
with the isometry factors $\sqrt m$ and $1/\sqrt m$ retained.
Its coefficient cover is $\mathbb C[z,z^{-1}]\to
\mathbb C[w,w^{-1}]$, $z\mapsto w^m$, with the full basis
$1,w,\ldots,w^{m-1}$ and the multiplication carry $z$.

For the same original remainder map $J_N$, let $C_N,C_{N,\theta}$
be the actual least-norm sections, and retain the literal relation
columns $B_{\rm rel}$. The complete relation identity is
\[
\begin{gathered}
 C_N-C_{N,\theta}=B_{\rm rel}X_{N,\theta},\qquad
 X_{N,\theta}=(B_{\rm rel}^*M_{N,\theta}B_{\rm rel})^{-1}
                  B_{\rm rel}^*M_{N,\theta}C_N,\\
 G_N=\overline G_N+\mathscr D_N,\qquad
 \mathscr D_N=\int(C_N-C_{N,\theta})^*M_{N,\theta}
                         (C_N-C_{N,\theta})\,\frac{d\theta}{2\pi}.
\end{gathered}
\]
Every relation has its original theta primitive. Empty relation
columns give $\mathscr D_N=0$. With the actual weighted source
constant $\kappa_{a,N}$ and $\delta_L=\kappa_{a,N}/(e^{aL}-1)<1$,
the full inverse-chord and block-matrix proof gives
\[
 (1-\delta_L^2)G_N\preceq\overline G_N\preceq G_N.
\]
For the four-volume comparison, compute the common bound
$\delta_L=\kappa_{a,N_*}/(e^{aL}-1)<1$ on the retained maximal
source $N_*=2q+2$ and restrict it to each of the four endpoint
spaces, as in (H42)--(H43). Its error is then at most
$2q\log(1/(1-\delta_L^2))$.  This bound is attached to the arithmetic
average of the canonical quotient forms.  The exact comparison
is \(G_N=\overline G_N+\mathscr D_N\), with the displayed
positive integral of original relation corrections; the
continuous equal-image complex and its cochain contraction
realizing \(\overline G_N\) are given in R64 and the full HCD
proof.  The individual phase-source comparisons have their
stronger current common-polynomial-source controls in
(HC13)--(HC14).  Their averaging retains this separate
positive correction and all of its original entries. The supplied two-point sharpness
is strengthened by the explicit compactly supported full-circle
source in (HDI22)--(HDI24), with original calibration mass seven
and phase factors $e^{\pm i\theta}$.


\paragraph{Exact graph-to-average continuation of R63.}
Retain the same $\mathcal P_N,E,J_N,B_N,M_N,M_{N,\theta}$ and
the original graph $\mathcal A_N=(\mathcal V_N,D_N)$, where
$D_N=\mathsf T_k\mathcal V_N-LJ_N$. The complete proof
(PGB.1)--(PGB.13), with the original all-vector source constants
and the circumference domain (PGB.8), gives the actual additional
positive form
\[
 \begin{gathered}
 H_N=B_N^*M_NB_N,\quad T_N=D_NB_N,\quad
 Q_{D,N}=(I+T_NH_N^{-1}T_N^*)^{-1},\\
 \mathcal E_N^{\rm aug}=C_N^*D_N^*Q_{D,N}D_NC_N>0,\qquad
 \widehat G_N-\overline G_N
     =\mathcal E_N^{\rm aug}+\mathscr D_N^{\rm hol}.
 \end{gathered}
\]
The relation-space inverse is the unique empty inverse when
$N=q-1$, so the composite through that space is zero. Here
$\mathscr D_N^{\rm hol}$ is exactly the already displayed
original relation gap $G_N-\overline G_N$. The proof computes
the actual defect map
\[
 x\longmapsto\left(Q_{D,N}^{1/2}D_NC_Nx,
       [\theta\mapsto(C_N-C_{N,\theta})x]\right)
\]
into the direct sum of the original defect Hilbert target and
the original relation-field space. Its Gram is the displayed
difference and its kernel is zero. The zero kernel is proved
by the actual leading coefficient: for $P$ of degree $n$ with
leading coefficient $p_n\ne0$, $\tau_kP(\Sigma)$ has $s_1$
degree $n+4m+1$ with leading coefficient $p_n/(4m+1)$, while
$F_{J_NP}$ has $s_1$ degree at most $4m-1$.
The full-unit source consequently makes $D_N$ injective.
Thus $\overline G_N\preceq G_N<\widehat G_N$ with a strict
finite graph cost. Repetition of the same polynomial coefficients
is the contracting cochain map inducing the identity
$(E,\widehat G_N)\to(E,\overline G_N)$; PGB.13 supplies its
complete isometric dilation.

For the four endpoints use one original circumference, chosen
from the common maximal source $N_*=2q+2$ whenever the inverse
formulas are used. PGB.26 gives the exact signed difference
as the four logarithmic determinants of
$\overline G_N^{-1/2}\widehat G_N\overline G_N^{-1/2}$,
with signs $(+,+,-,-)$ at $(q-1,q,2q-1,2q)$.
The original averaged-quotient bound above stays attached to
$\overline G_N$ and its original relation gap. The new signed
graph-to-average expression retains every extra term; positivity
of each individual logarithm does not give a sign for that
four-term expression.

\paragraph{Result R64: continuous cochain descent and the exact completion map.}
The full proof (HCD1)--(HCD29), including the actual source bounds
(HCD2a)--(HCD2c), considers relation fields $\mathscr B$ and the
common-image fields $\mathscr P^{\rm eq}$ over $d\theta/(2\pi)$.
At the fixed cutoff $N$, choose $0<\eta<1$ and retain the actual
weighted source constant $\kappa_{\beta,N}$. The explicit
circumference
$L\ge\beta^{-1}\log(1+\kappa_{\beta,N}/\eta)$ gives the uniform
equivalence in (HCD2a)--(HCD2c) used in the following bounded maps.
Repetition and Bochner mean give the bounded coefficient isomorphism
\[
 [\mathscr B\hookrightarrow\mathscr P^{\rm eq}]
 \simeq[\chi\mathcal P_{N-q}\hookrightarrow\mathcal P_N]
       \oplus[\mathscr B^0\xrightarrow{I}\mathscr B^0],
 \qquad hP=P-\Delta\operatorname{av}P .
\]
The full calculation proves $dh+hd=I-\Delta\operatorname{av}$,
and the quotient is the original $E=\mathbb C[S]/(\chi)$ with
metric $\overline G_N$. Its relation gap retains the weighted
cross term. If $\overline C_N=\int C_{N,\theta}\,d\theta/(2\pi)$,
$a_C=C_N-\overline C_N$ and $h_\theta=C_{N,\theta}-\overline C_N$,
then
\[
 \int a_C^*M_{N,\theta}h_\theta\,\frac{d\theta}{2\pi}=a_C^*M_Na_C,
 \qquad
 \mathscr D_N=\int h_\theta^*M_{N,\theta}h_\theta
                    \,\frac{d\theta}{2\pi}-a_C^*M_Na_C .
\]

For every $L>0$, the full shifted-density proofs (CQ1)--(CQ29) and
(HDI11)--(HDI21) identify the subsequent completed phase quotient with the weighted
values at positive-mass sampled roots. Its exact map from $E$ has
kernel $(p_\theta)/(\chi)$, where $p_\theta$ is the squarefree
sampled product and $\chi$ retains every primary multiplicity.
The original one-fold packet satisfies $(\chi_{h,1})=(h)$; its
selected-root mass is
$|g^{(r)}(\rho)/h^{(r)}(\rho)|^2/L>0$, with the original leading
coefficient retained. The complete primary-coordinate inverse gives
the nilpotent kernel at every sampled root and the full block at
every unsampled root.

The diagonal root projections give an explicit measurable field.
Only finitely many phases can sample a fixed root of $\chi$, so
its completed direct integral is zero. The evaluation maps form
the actual cochain map from the continuous common-image complex
to this completed complex; its map on degree-one cohomology has
kernel all of $E$. On an original split label $\lambda$, its
inverse image of the receiving represented zero is the full
represented coefficient kernel, while the inverse image of $\tau$
is the singleton $\{\tau\}$. This specifies the operation and its
kernel with the original support and all finite classes present.


\paragraph{The corrected-graph completion and its exact map in R64.}
The full proof (PGB.14)--(PGB.21b) retains the original scalar
graph density $r_k$, full-unit isometry $\mathscr J_k$,
mixed observation $Y=(0,L)$, its actual adjoint $g=\mathscr J_k^*Y$,
and $Y_\perp=(I-\mathscr J_k\mathscr J_k^*)Y$.
The density proof uses the original exponential moments and
proves the onto multiplication isometry for the literal $\chi$;
it does not presume polynomial density. The completed source,
its inverse coordinate map, and its completed relation space are
\[
 \begin{gathered}
 \mathscr K_\infty=
 \mathscr J_kL^2(r_kdu)\mathbin{\oplus^\perp}Y_\perp E,
 \quad\Omega_k=Y_\perp^*Y_\perp>0,\\
 \Xi(f,x)=\mathscr J_kf-Y_\perp x,\qquad
 \Xi^{-1}(v)=
 (\mathscr J_k^*v,-\Omega_k^{-1}Y_\perp^*v),\\
 \overline{\mathcal A(\chi\mathbb C[S])}
       =\mathscr J_kL^2(r_kdu),\qquad
 \mathscr K_\infty/\overline{\mathcal A(\chi\mathbb C[S])}
       \simeq(E,\Omega_k).
 \end{gathered}
\]
The quotient map is the displayed second coordinate, so its
kernel is precisely that closed relation subspace. Its vertical
part in the original ordered two-component target is exactly
$\{(0,-Lx):x\in E\}$. This is a closed linear relation;
the full proof constructs original polynomials with fixed
nonzero remainder, observation tending to zero and defect
tending to $-Lx$, proving that the original defect operator
on the polynomial observation source is not closable.

The first-coordinate projection has the exact formula
$\pi_1\Xi(f,x)=\mathscr J_k^0I_{r,m}(f+g_x)$, where
$\mathscr J_k^0f=\prod_i v_h(s_i)f(u)$ and
$I_{r,m}:L^2(r_kdu)\to L^2(m_kdu)$ is the unchanged-function
contraction. It defines the bounded cochain map
\[
 [\mathscr J_kL^2(r_kdu)\hookrightarrow\mathscr K_\infty]
 \longrightarrow
 [\mathscr J_k^0L^2(m_kdu)\xrightarrow{I}
                         \mathscr J_k^0L^2(m_kdu)].
\]
On original polynomials it is literally
$\mathcal AP\mapsto\mathcal VP$. Its degree-one map is
$(E,\Omega_k)\to0$, with kernel all of $E$, matching the
original unweighted completion. The finite contraction of R63
and this completed map have their distinct, fully specified
domains and kernels. Applying a fixed retained support label
preserves the represented kernel; external absence retains
its singleton preimage.


\paragraph{The strict mixed source inside the same completed quotient.}
The complete original-source calculation (PGMB.1)--(PGMB.8)
now computes the interior of $\Omega_k$. With the unchanged
positive scalar $1+|\tau_k|^2$ define
\[
 \begin{gathered}
 L_{\rm low}x=
 \frac{(\prod_i v_h(s_i))F_x(s)}{\sqrt{1+|\tau_k|^2}},
 \quad L_{\rm low}:E\longrightarrow\mathscr H_k,\\
 \mathscr D_k^{\rm mix}x=
 \frac{\overline{\tau_k}F_x(s)}{1+|\tau_k|^2}-g_x(u),
 \quad \mathscr D_k^{\rm mix}:E\longrightarrow
 L^2\bigl(\mathbb R^k,(1+|\tau_k|^2)\prod_iw(t_i)\,dt_i\bigr).
 \end{gathered}
\]
The first original product integral has Gram $K_{\rm low}$;
the orthogonal sum-coordinate projection in PGM.3--5 gives
$(\mathscr D_k^{\rm mix})^*\mathscr D_k^{\rm mix}
=\Omega_k-K_{\rm low}$. Thus the specified map
$\mathfrak U_kx=(L_{\rm low}x,\mathscr D_k^{\rm mix}x)$
is an isometry from $(E,\Omega_k)$ into that original direct
sum. Its complete left inverse is
$\Omega_k^{-1}(L_{\rm low}^*a+
(\mathscr D_k^{\rm mix})^*b)$; the reverse composite is the
orthogonal projection onto $\mathfrak U_k(E)$.
Writing the already computed quotient as
$j_\infty v=-\Omega_k^{-1}Y_\perp^*v$, its composition is
\[
 \begin{gathered}
 \mathfrak C_k=\mathfrak U_kj_\infty,\qquad
 \mathfrak C_k\Xi(f,x)=(L_{\rm low}x,\mathscr D_k^{\rm mix}x),\\
 \ker\mathfrak C_k=\mathscr J_kL^2(r_kdu),\qquad
 \|\mathfrak C_k\Xi(f,x)\|^2=x^*\Omega_kx.
 \end{gathered}
\]
This is the actual quotient contraction and its whole relation
kernel, since $\mathfrak U_k$ is injective and $j_\infty$
has the displayed kernel. The complete fixed-sum polynomial
argument PGM.6--9 proves
$\ker\mathscr D_k^{\rm mix}=0$ for every $k\geq2$ and
$\mathscr D_1^{\rm mix}=0$. Consequently for $k\geq2$
\[
 (\Omega_k-K_{\rm low})^{-1}
 (\mathscr D_k^{\rm mix})^*\operatorname{pr}_2\mathfrak C_k
 =j_\infty,
\]
by the computed mixed Gram. For $k=1$, the second component
vanishes and the positive first Gram retains the full quotient.
All maps act coordinatewise on each retained split label;
their represented-zero kernels are the products of the exact
displayed coefficient kernels. The empty present tuple and
the singleton inverse image of external $\tau$ are unchanged.
Composition from a larger or proper-source realization retains
the whole earlier realization kernel specified in PGM.10;
this calculation asserts no descent through another quotient.

\paragraph{Result R65: quadratic generator transfer with an exact residue norm.}
The complete proof (HG1)--(HG25), including (HG11a), uses the
actual relation-cost endomorphism $\mathsf D=G_N^{-1}\mathscr D_N$,
the original $T_t=A+tR-kI/2$, and $R=e_0\ell$. It proves
\[
\begin{gathered}
 |\epsilon_{\overline G_N}(t)-\epsilon_{G_N}(t)|
 \le\frac{\|[\mathsf D,A]+t[\mathsf D,R]\|_{G_N}}{1-\delta_L^2}
 \le\frac{\delta_L^2}{1-\delta_L^2}\|T_t\|_{G_N}.
\end{gathered}
\]
The exact square-root isometry and Sylvester integral preserve
the full Hermitian error. The sharper scalar constant follows
from the complete spectral-projection decomposition of the
original self-adjoint discrepancy: its commutator norm is at
most its spectral diameter times the original generator norm.

For any actual self-adjoint discrepancy $B=G^{-1}(H-G)$, let
$x=e_0$, $y=G^{-1}\ell^*$ and retain
\[
\begin{gathered}
 a_0=\|Bx\|_G^2,\ c_0=\|x\|_G^2,\ d_0=\|y\|_G^2,
 \ f_0=\|By\|_G^2,\quad
 p_0=\langle x,Bx\rangle_G,\ r_0=\langle y,By\rangle_G,\\
 T_{\rm res}=a_0d_0+c_0f_0-2p_0r_0,\qquad
 \Delta_{\rm res}=(a_0c_0-p_0^2)(d_0f_0-r_0^2),\\
 \|[B,R]\|_G^2=
 \frac{T_{\rm res}+\sqrt{T_{\rm res}^2-4\Delta_{\rm res}}}{2}.
\end{gathered}
\]
The full coordinating residue proof is reproduced in (HG12)--(HG15),
including both negative cross terms, singular Gram matrices,
rank zero, rank one, repeated singular values and $q=1$.
The primary quantity remains $\|[B,A]+t[B,R]\|_G$ with its
complete cancellation.

Equations (HG16)--(HG20) give explicit $L,s,m,J$ from the actual
weighted and derivative source Grams. The holomorphic Schur
inverse is certified on the specified smaller strip, with its
neighborhood, and the finite mean obeys
\[
 \alpha G_N\preceq H_{m,J}\preceq\beta G_N,\qquad
 \alpha=(1-\xi_J)(1-\delta_L^2-\varepsilon)>0,
 \quad\beta=1+\varepsilon.
\]
The exact discrepancy is $B_{m,J}=-\mathsf D+\mathsf E_{m,J}$,
where the finite-tail spectrum lies in $[-\rho,\varepsilon]$,
$\rho=\xi_J+(1-\xi_J)\varepsilon$. The proof retains its entire
commutator and gives the corresponding spectral-diameter bounds.
The finite representative retains both ordered theta boundary
terms in (HG23). The resulting original-metric radius enters the
complete Laplacian numerical-range parabola (HG25); the same proved
metric comparison gives the residue-coefficient bounds (HG24).


\paragraph{Exact discrepancy inputs supplied to R65 by the graph continuation.}
The original estimates above use
$\mathsf D=G_N^{-1}\mathscr D_N^{\rm hol}$ and retain their
specified source and circumference constants. The exact graph
continuation computes, rather than omits, the further form
$\mathcal E_N^{\rm aug}$ in PGB.6. Relative to the original
metric, the graph discrepancy is the actual self-adjoint
endomorphism
\[
 B_N^{\rm aug}=G_N^{-1}(\widehat G_N-G_N)
             =G_N^{-1}\mathcal E_N^{\rm aug}.
\]
It is positive because
$\langle x,B_N^{\rm aug}x\rangle_{G_N}
=x^*\mathcal E_N^{\rm aug}x>0$ for $x\ne0$;
$G_NB_N^{\rm aug}=\mathcal E_N^{\rm aug}$ is Hermitian,
which proves self-adjointness in precisely that original metric.
For the comparison based at the averaged metric the endomorphism
is instead the explicitly typed
\[
 B_N^{\rm graph,av}
   =\overline G_N^{-1}(\widehat G_N-\overline G_N)
   =\overline G_N^{-1}
        (\mathcal E_N^{\rm aug}+\mathscr D_N^{\rm hol}).
\]
Multiplication by $\overline G_N$ proves its self-adjointness
and positivity in $\overline G_N$. Thus the generator terms
for these comparisons are respectively
$[B_N^{\rm aug},A]+t[B_N^{\rm aug},R]$ and
$[B_N^{\rm graph,av},A]+t[B_N^{\rm graph,av},R]$,
with the unchanged $A,R,T_t$. This identifies all entries
needed by the existing transfer calculation. The present
continuation supplies no new small bound for these additional
commutators; the existing bound on $\mathsf D$ remains a bound
for its stated original term.


\paragraph{The calculated finite and completed generator controls in R65.}
The complete proof PGGT.1--8 now evaluates the transfer for
three actual rows: $(G,H)=(G_N,\widehat G_N)$,
$(\overline G_N,\widehat G_N)$ and $(K_{\rm low},\Omega_k)$.
Their respective discrepancies are exactly
\[
 B=G_N^{-1}\mathcal E_N^{\rm aug},\qquad
 B=\overline G_N^{-1}
           (\mathcal E_N^{\rm aug}+\mathscr D_N^{\rm hol}),\qquad
 B=K_{\rm low}^{-1}(\Omega_k-K_{\rm low}).
\]
Each is self-adjoint and nonnegative in its displayed original
$G$ metric, by $GB=H-G$ and the preceding positive Gram
proofs. Keep its complete original spectral decomposition
$B=\sum_i b_iP_i$ with all repeated eigenspaces, and set
$s_i=\sqrt{1+b_i}$ and $\mathsf S=\sum_i s_iP_i$.
Multiplying proves $\mathsf S^*G\mathsf S=H$ and
$\mathsf S^{-1}=\sum_i s_i^{-1}P_i$. Thus with the original
$T_t=A+tR-kI/2$ the full error is
\[
 \begin{gathered}
 Z=\sum_{i,j}\frac{s_i-s_j}{s_j}P_iT_tP_j,\qquad
 \mathsf S X_H\mathsf S^{-1}-X_G=Z+Z^{\dagger_G},\\
 |\epsilon_H(t)-\epsilon_G(t)|
 \leq\|Z+Z^{\dagger_G}\|_G
 \leq\frac{\|[B,A]+t[B,R]\|_G}{1+b_1}
 \leq\frac{b_r-b_1}{1+b_1}\|T_t\|_G.
 \end{gathered}
\]
The displayed denominator is the actual spectral minimum.
PGGT.4--5 proves the entire identity by ordered multiplication,
then proves the middle estimate by integrating
$e^{-u\mathsf S}[B,T_t]e^{-u\mathsf S}$ with its exact
$1/(2\sqrt{1+b_1})$ factor. The full spectral projection
decomposition supplies the last bound; the residue term remains
$[B,R]v=(Be_0)\langle G^{-1}\ell^*,v\rangle_G
-e_0\langle BG^{-1}\ell^*,v\rangle_G$, including the
negative term and its full cancellation with $[B,A]$.

For the third row the spectral endpoints are precisely
$b_1=\lambda_{k,1}$ and $b_r=\lambda_{k,q_k}$ of PGMB.9.
At $k=1$ their common zero gives exact equality of the
two Hermitian generators. The original Hilbert-source map
$\mathfrak U_k=(L_{\rm low},\mathscr D_k^{\rm mix})$
realizes the completed generator by
\[
 \begin{gathered}
 \widetilde T_t=\mathfrak U_kT_t\Omega_k^{-1}\mathfrak U_k^*,
 \quad\widetilde T_t\mathfrak U_k=\mathfrak U_kT_t,\\
 \ker\widetilde T_t
 =\mathfrak U_k(\ker T_t)\mathbin{\oplus^\perp}
                                  \mathfrak U_k(E)^\perp,
 \qquad
 \|\widetilde T_t+\widetilde T_t^*\|
                           =\epsilon_{\Omega_k}(t).
 \end{gathered}
\]
PGGT.7--8 proves these on the original target by decomposing
each vector into $\mathfrak U_kx+v_\perp$ and using
$\mathfrak U_k^*\mathfrak U_k=\Omega_k$. Thus the new
estimate controls that actual realized Hermitian generator;
its constants retain the entire mixed spectrum and make no
unproved claim of smallness as $k$ grows.

\paragraph{Result R66: the original degree cost and full continuous closure.}
The complete proof (HD1)--(HD10) fixes $h,k,L$ and a real sample
$u_0$, with $S_0=k/2+iu_0$ and $a_N=(1,S_0,\ldots,S_0^N)$.
For the unchanged original polynomial Gram it constructs the
polynomial $P_N$ with its displayed coefficient vector $c_N$:
\[
 \Lambda_N=a_NM_N^{-1}a_N^*\uparrow\infty,\qquad
 c_N=M_N^{-1}a_N^*/\Lambda_N,
 \qquad P_N(S_0)=1,\quad\|\mathcal V P_N\|^2=1/\Lambda_N.
\]
The divergence is proved by the original exponential moments,
Hilbert representation, holomorphic Fourier uniqueness and the
Gaussian measure argument. At an actually positive sampled mass
$\omega_0=(2\pi/L)m_{h,k}(u_0)$ this gives the exact degree cost
\[
 \kappa_{a,N}\ge(e^{aL}-1)(\omega_0\Lambda_N-1)
       \longrightarrow\infty .
\]
The same columns have full holonomy mean squared norm tending to zero,
while their selected phase squared norm remains at least $\omega_0$.
Evaluation factors through the original $E$ exactly when
$\chi(S_0)=0$, with the fully calculated primary kernel.

The coordinating complete proof (CR1)--(CR14), including (CR12a)--(CR12c),
extends this to the entire original source. Multiplication by the
literal $\chi$ is the onto isometry
\[
 L^2(|\chi|^2m\,du)\xrightarrow{\ \times\chi\ }L^2(m\,du),
 \qquad f\longmapsto\chi f,
\]
with its division inverse defined almost everywhere. Polynomial
density therefore proves that the continuous relation closure is
the whole Hilbert source. If $\Pi_N$ is its actual finite relation
projection and $s:E\to\mathcal P_{q-1}$ its fixed remainder section,
then
\[
 R_N=(I-\Pi_N)s,\qquad J_NR_N=I_E,\qquad
 G_N=R_N^*R_N\downarrow0,
 \qquad\overline{\operatorname{graph}J}=L^2(m\,du)\times E.
\]
All limits here fix the original $h,k$ and retain the complete
finite coefficient frame. At singular phase Grams the exact
receiving space is $T_\theta(\mathcal P_N)/T_\theta(\mathcal B_N)$,
with kernel $J_N(\ker T_\theta)$ on $E$. Its measurable minimum
form gives $0\preceq\overline G_N\preceq G_N$ by testing the
original $R_N$. The independent full completion proof further
gives the relation gap as the integral of the squared projection
of $T_\theta R_N$ onto $T_\theta\mathcal B_N$. Every finite
$\overline G_N$ is positive for each original $k\ge1$: outside
a countable set of one-fold phases all sampled weights are
positive. At an actual positive sampled root, its fixed-phase
quotient/source ratio diverges by (CR14). These exact maps and
limits impose no unstated rate in the distinct tensor parameter
$k$ and do not identify the relative gap with its absolute limit.

\paragraph{Exact simultaneous degree limits and morphisms extending R66.}
The complete affine-distance formula in PGB.22 is
\[
 x^*\widehat G_Nx=x^*\Omega_kx+
 \inf_{Q\in\mathcal P_{N-q}}
 \|\iota(sx+\chi Q)-g_x\|_{L^2(r_kdu)}^2,
 \qquad (\iota P)(u)=P(k/2+iu).
\]
The original weighted relation density, proved in PGB.15--17,
therefore gives at fixed original $h,k$
\[
 \widehat G_N\downarrow\Omega_k>0,\qquad
 G_{W,N}\downarrow0,\qquad\Delta_N\downarrow\Omega_k,
 \qquad G_N\downarrow0,\qquad\overline G_N\to0.
\]
The last two limits are the original unweighted limits already
stated above. For every circumference the phase quotient is
defined by minimum source images, including singular finite
Grams, and its projected relation gap proves
$0\preceq\overline G_N\preceq G_N$ without inverting a null
coordinate. The full proof also establishes strict positivity
of every finite $\overline G_N$ from the original sampled
weights. Consequently the graph augmentation tends to
$\Omega_k$ and the exact comparison obstruction is
\[
 \widehat G_N-\overline G_N\longrightarrow\Omega_k,
 \qquad
 \inf_{x\ne0}\frac{x^*\widehat G_Nx}{x^*\overline G_Nx}
 \geq\frac{\lambda_{\min}(\Omega_k)}{\|G_N\|}
                          \longrightarrow\infty.
\]
The inequality follows by bounding numerator and denominator
in the same original remainder frame. It excludes a constant
upper comparison independent of the polynomial cutoff;
the exact contracting identity in the other direction remains
PGB.12--13. At completion, the bounded first-coordinate
cochain map in PGB.21b sends $(E,\Omega_k)$ to the original
zero unweighted quotient with full coefficient kernel. None
of these fixed-$h,k$ degree limits asserts a rate in the
separate tensor parameter or deletes any retained primary
coordinate before its specified quotient map.


\paragraph{Finite tail costs and strict mixed persistence in R66.}
For the same actual graph minimum $\widehat R_N$ and the same
adjoint $g_x$, set
$f_Nx=\iota\widehat R_Nx-g_x\in L^2(r_kdu)$.
Substitution in PGB.19 and orthogonality prove the full identity
\[
 \widehat G_N
 =K_{\rm low}+(\mathscr D_k^{\rm mix})^*
            \mathscr D_k^{\rm mix}+f_N^*f_N
 =\Omega_k+f_N^*f_N.
\]
Its actual isometry is
$x\mapsto(L_{\rm low}x,\mathscr D_k^{\rm mix}x,f_Nx)$,
with all three original Hilbert targets as in PGMB.5--6.
The preceding degree limit gives $\|f_N\|\to0$ at fixed
$h,k$. For every $k\geq2$, PGM proves the strictly positive
mixed Gram $\Omega_k-K_{\rm low}$, so this entire mixed
part survives in the quotient while the third component
tends to zero. At $k=1$ its vanishing is the exact PGM identity.

Let $\lambda_{k,j}$ be the full generalized eigenvalues of
$(\Omega_k-K_{\rm low},K_{\rm low})$ in the retained
coefficient frame, and let $\nu_{k,N,j}$ be the eigenvalues
of $\Omega_k^{-1/2}f_N^*f_N\Omega_k^{-1/2}$. The complete
determinant calculation PGMB.9--10 yields
\[
 \log\frac{\det\widehat G_N}{\det K_{\rm low}}
 =\sum_{j=1}^{q_k}\log(1+\lambda_{k,j})
    +\sum_{j=1}^{q_k}\log(1+\nu_{k,N,j}),
 \qquad \nu_{k,N,j}\longrightarrow0.
\]
Every $\lambda_{k,j}$ is positive for $k\geq2$ and zero
for $k=1$. The exact extremal comparison constants are
$1+\lambda_{k,1}$ and $1+\lambda_{k,q_k}$, as proved by
testing the corresponding original-frame generalized
eigenvectors. The four unchanged endpoints cancel the entire
first sum with its literal coefficients $1+1-1-1=0$.
Their graph-volume expression is therefore the full signed
sum of the four tail costs, with $\Omega_k$ still present
in every tail congruence. These fixed-packet limits do not
assign any rate when $k$ and the endpoints grow together.

